% Methodology section (draft generated from the implemented code; keep in sync with configs/base.yaml).
% Symbols are chosen so that no letter has two meanings:
%   FiLM scale/shift: \phi, \psi     gate: g, \lambda     CB loss: \beta_{cb}     focal exponent: \kappa
\section{Methodology}

\subsection{Problem formulation}
For a field $i$ we observe a Sentinel-2 time series up to a cutoff day $\tau$,
$X_i^{\le\tau} = \{(\mathbf{x}_{t}, d_{t})\}_{t=1}^{T_i}$, where $\mathbf{x}_t \in \mathbb{R}^{11}$ holds ten
surface-reflectance bands (B2--B8, B8A, B11, B12) and the cloud probability, and $d_t$ is the acquisition day
relative to 1~January of the target year (negative for the previous year). We also observe static field
attributes $\mathbf{s}_i$ (elevation, slope, irrigation system, municipality) and the previous-year crop
$y_i^{\text{prev}} \in \{1,\dots,C\} \cup \{\varnothing\}$. The task is to predict the current crop
$y_i \in \{1,\dots,C\}$, $C = 18$, for any $\tau$; we report $\tau \in \{60, 152, 244, 335\}$
(1~March, 1~June, 1~September, 1~December).

\subsection{Temporal branch}
Each observation is mapped to $\mathbf{z}_t = [\tilde{\mathbf{x}}_t ; d_t/365] \in \mathbb{R}^{12}$, where
$\tilde{\mathbf{x}}_t$ is standardised with training-set statistics. A linear projection and a sinusoidal
encoding of the acquisition day give
\begin{equation}
  \mathbf{h}^{(0)}_t = \mathrm{ReLU}\big(W_{\text{in}}\,\mathrm{LN}(\mathbf{z}_t)\big) + \mathrm{PE}(d_t),
\end{equation}
which a five-block Transformer encoder (model width 64, two heads) turns into tokens
$\mathbf{H} = [\mathbf{h}_1, \dots, \mathbf{h}_T]$, with padded steps masked out of attention.

\subsection{Static branch}
Continuous attributes are standardised; categorical attributes are embedded. An MLP produces a static
summary $\mathbf{h}_s \in \mathbb{R}^{64}$, and each attribute is also projected to its own token, giving
$\mathbf{S} \in \mathbb{R}^{4 \times 64}$.

\subsection{Fusion}
Temporal tokens attend to the static tokens,
\begin{equation}
  \mathbf{h}'_t = \mathrm{LN}\big(\mathbf{h}_t + \mathrm{MHA}(\mathbf{h}_t, \mathbf{S}, \mathbf{S})\big),
  \qquad
  \mathbf{h}_{\text{temp}} = \max_{t \in \mathcal{V}} \mathbf{h}'_t ,
\end{equation}
where $\mathcal{V}$ is the set of valid (non-padded) steps. Feature-wise linear modulation (FiLM) then lets
the static summary scale and shift the temporal representation:
\begin{equation}
  \mathbf{h}_{\text{FiLM}} = \big(\mathbf{1} + \phi(\mathbf{h}_s)\big) \odot \mathbf{h}_{\text{temp}} + \psi(\mathbf{h}_s),
\end{equation}
with $\phi$ and $\psi$ linear maps initialised to zero, so training starts from the unmodulated features.

\subsection{Rotation-aware adaptive gate}
Let $\boldsymbol{\mu}_c(b)$ be the mean standardised band vector of class $c$ in five-day bin $b$ of the
year, computed on the training fields. For a field with known previous crop $p$, the divergence of the
current season from that crop's expected development is
\begin{equation}
  D_{\text{rot}} = \frac{1}{|\mathcal{C}|} \sum_{t \in \mathcal{C}}
  \Big(1 - \cos\big(\tilde{\mathbf{x}}^{\text{band}}_t, \boldsymbol{\mu}_p(\lfloor d_t / 5 \rfloor)\big)\Big),
\end{equation}
where $\mathcal{C}$ are the current-season observations up to $\tau$. The gate is
\begin{equation}
  g = \sigma\big(\mathbf{w}_g^{\top}[\mathbf{h}_{\text{temp}} ; \mathbf{e}_p ; \mathbb{1}_{|\mathcal{C}|>0}] + b_g
      - \lambda\, D_{\text{rot}}\big), \qquad \lambda = \mathrm{softplus}(\lambda_0) > 0,
\end{equation}
with $\mathbf{e}_p$ a learned embedding of the previous crop and $g = 0$ when the previous crop is unknown.
The fused representation is
\begin{equation}
  \mathbf{h} = \mathbf{h}_{\text{FiLM}} + g\, \mathbf{e}_p ,
\end{equation}
and an MLP head outputs class probabilities $\mathbf{p} = \mathrm{softmax}(f(\mathbf{h}))$.

\subsection{Class-balanced focal loss}
With $n_c$ training fields in class $c$, the effective-number weight is
$\alpha_c \propto (1 - \beta_{cb}) / (1 - \beta_{cb}^{\,n_c})$, normalised to sum to $C$, and
\begin{equation}
  \mathcal{L} = - \alpha_{y}\,(1 - p_{y})^{\kappa} \log p_{y},
\end{equation}
with $\beta_{cb} = 0.999$ and $\kappa = 1$. (A stronger setting, $\beta_{cb} = 0.9999$ and $\kappa = 2$, weighted the
rarest class about 240 times the most common one and led one training run to ignore the satellite series and
repeat the previous-year crop; the milder setting avoids this.)

\subsection{Training protocol}
All runs follow the training setup of the baseline's official implementation, so that run A is a faithful
reproduction and runs B--D are compared with it under identical conditions. Models are trained on the 2022
fields (stratified 90/10 train/validation split) with a cutoff day drawn uniformly from the year for each sample,
so a single model serves every date; for a cutoff $\tau$ the input holds the observations of the 365 days before
$\tau$. Optimisation uses Adam (learning rate $10^{-3}$, weight decay $5\times10^{-8}$, batch size 512, constant
learning rate) for at most 120 epochs, with early stopping (patience 10) on the mean macro-F1 over the four
benchmark dates on the 2022 validation split. The 2023 fields are used only for testing, never for model
selection. Each configuration is trained with three random seeds and we report the mean and standard deviation.

\subsection{Ablation design}
Run A is the baseline of Gen\'e-Mola et al.\ (2026), re-implemented from its official code: a single Transformer
encoder (linear $12\to64$, ReLU, five post-norm encoder layers with two heads and feed-forward width 128,
dropout 0.018, max-pooling, ReLU, linear classifier) without positional encoding, in which each ancillary value
(previous-year crop, irrigation system, municipality, elevation, slope; scaled to $[-1,1]$) is appended as one
additional time step. The only remaining difference is that the 365-day window is given as the list of observed
dates with a padding mask rather than as a zero-filled daily grid. Run B replaces this model with the dual-branch
model with cross-attention and FiLM (previous crop always added, $g = 1$). Run C adds the rotation-aware gate,
and run D adds the class-balanced focal loss. From B to D each run differs from the previous one in exactly one
component.
